\subsubsection{非单值超势与 $q$-分次}\label{nonsingle}

在迄今所给出的框架中，表示辫群在 $\H$ 上作用的矩阵其元素取整数值，因为相对同调具有整结构。现在我们要考虑 $\W$ 非单值的情形。更准确地说，我们将考虑像 Chern-Simons 泛函 $\W$ 那样的全纯函数：其实部单值，但虚部只模 $2 \pi \Z$ 有定义。只要把 $\U$ 换成使 $\W$ 单值的最小覆盖 $\hat \U$，就可以化归到我们迄今所用的框架。

转到 $\hat\U$ 的代价是临界点的数目变成无穷多，因为 $\U$ 中每个临界点在 $\hat \U$ 中都有无穷多个原像。$\hat \U$ 中的一个临界点 $\hat I$ 等同于 $\U$ 中的一个临界点 $I$ 连同在 $I$ 处对 $\W$ 的一条分支的选取。局部地，我们可以任取临界点 $I$ 的一个原像 $\hat I_0$，并把 $\W$ 的值比 $\hat I_0$ 处的值平移 $-2 \pi i n$ 的那个临界点记作 $\hat I_n$。

Stokes 因子与传输矩阵现在是无穷维矩阵，但其矩阵元可以用我们已描述的技术计算，且取整数值。由于这些计算只依赖于 $\W$ 的梯度，这些矩阵与覆叠变换 $\hat I_n \to \hat I_{n+1}$ 对易。引入一个取值于 $\C^*$ 的变量 $q$ 是方便的。于是，若我们简单地把 $q^n\psi_I$ 写成 $\psi_{\hat I_n}$ 的符号缩写，就可以把元素为整数的无穷维矩阵替换为元素为整系数 $q$ 的 Laurent 多项式的有限维矩阵。我们把这样的矩阵记作 $S[\ell_I^J;q]$。为了追踪在 $\ZZ$ 中移动时 $\W$ 的提升，编码临界点置换的形式单值矩阵以及与梯度流相联系的 Stokes 矩阵都取值于 $q$ 的幂，以追踪沿路径 $\W$ 的变化。传输矩阵的矩阵元现在通过对混合路径（远离 Stokes 墙处的连续演化与穿过 Stokes 墙的梯度流）求和得到，权重为 $\pm q^n$，其中 $-2\pi i n$ 是超势沿路径的变化，而符号一如既往地涉及针的定向。

对这一切还有一种等价的描述。由于 $\U$ 容许非单值超势 $\W$，其第一 Betti 数为正，人们可以在超对称量子力学中引入一个“theta 角” $\theta$，对 $\W$ 跳跃 $-2\pi i n$ 的路径加权 $e^{in\theta}$。其效果是把 $\U$ 的相对同调换成一个扭曲的相对同调，取值于单值性为 $q=e^{i\theta}$ 的平坦线丛。然后我们考虑由这个平坦丛扭曲的同调的 Gauss-Manin 联络。该联络的和乐矩阵其元素为整系数 $q$ 的 Laurent 多项式，并且可以按刚才描述的方式计算。
% [Tutor Note: 易错点] 本节的 $q$ 是形式变量（用于把无穷维整矩阵压缩为有限维 Laurent 多项式），而第~\ref{clco}~节之后（如 (\ref{torg})、(\ref{zorg})）的 $q$ 是取数值的 Jones 参数 $q=\exp(2\pi i/(\k^\vee+2))=\exp(-2\pi i/b^2)$。两者在本文中同名，但一为形式、一为数值，读到 $q^n$、Laurent 多项式时不要误代入数值。




\subsection{四维解计数的经典描述}\label{clco}

现在让我们考虑真正感兴趣的问题——$\R\times C\times \R_+$ 上超对称方程 (\ref{bpseqns}) 的依赖时间的解。解之所以依赖时间，是因为边界条件依赖时间——我们允许 $\Mthree=C\times \R_+$ 边界上的奇异磁单极子的位置 $z_i$ 随时间变化。虽然我们通常假定磁单极子位置作绝热演化，但四维解的计数是拓扑的，绝热假定并非必需。



在最简单的设定中，边界处的奇异磁单极子股在时间中编织，超势全纯地依赖于某些参数（即各股在 $C$ 中的位置），而这些参数随时间演化。概略地说，方程取“受迫梯度流”的形式
\begin{equation} \label{flow}
\frac{\d \bar w_{\bar i}}{\d t} = -e^{i \alpha} g^{\bar i j}  \frac{\partial \W(w,z(t))}{\partial w_j}
\end{equation}

我们假定这些奇异磁单极子在 $t=-\infty$ 附近从位置 $\vec z_i=(z_{1},\dots,z_{k})$ 出发，并在 $t=+\infty$ 附近终止于位置 $\vec z_f=(z'_{1},\dots,z'_{k})$。事实上，我们可以假定奇异磁单极子的位置 $z_i(t)$ 除在某个有界区间 $-T<t<T$（$T$ 为某值）内之外都是常数，而在这个区间内它们沿其参数空间 $\ZZ$ 中的一条路径 $\P$ 运动。在这种情形下，我们可以寻找受迫梯度流方程的解：它们在 $\W(w,\vec z_i)$ 的某个指定临界点 $I$ 出发，终止于 $\W(w,\vec z_f)$ 的某个指定临界点 $J$。这类解的“数目”（每个解按费米子行列式的符号加权）是一个拓扑不变量——在路径 $\P$ 或 $\U$ 上度规的形变下不变。我们把这个不变量记作 $\frak N^J_I$。（由于一个马上就会明白的原因，我们实际上只想在 $\vec z_i$ 与 $\vec z_f$ 不在 Stokes 墙上时定义 $\frak N^J_I$。）

在第~\ref{time}~节中，我们已经为这种情形联系了一个整数不变量 $\frak n^J_I[\P]$。$\frak n^J_I[\P]$ 是沿从 $\vec z_i$ 到 $\vec z_f$ 的路径 $\P$ 作平行移动的 Gauss-Manin 联络的矩阵元。可以把 $\frak n^J_I[\P]$ 看作这样的展开系数：把相对同调 $H(\U,\U_<)_{\vec z_i}$ 中的一个针 $\J_I$ 沿路径 $\P$ 用 Gauss-Manin 联络搬运，然后再用构成 $H(\U,\U_<)_{\vec z_f}$ 的一组基的针 $\J'_J$ 来表示：
\be\label{zelx}\J_I=\sum_J\frak n^J_I \J'_J.\ee
我们断言，事实上
\be\label{elx}\frak N^J_I=\frak n^J_I.\ee
这一关系的重要性在于：$\frak N^J_I$ 才是我们想要的东西，即依赖时间的解的计数；而 $\frak n^J_I$ 更容易计算，因为它只需研究与时间无关的问题。
% [Tutor Note: 物理图像] (\ref{elx}) 是本节的核心：左边是“真正”的四维（依赖时间）解数，右边是纯三维（与时间无关）的 Gauss-Manin 搬运系数；二者相等意味着时间依赖只把针连续地搬运，不会额外产生或消灭解。实现这一等价的正是“强迫梯度流把 $\J_I$ 连续形变为 $\tilde\J_I$、而相对同调类是离散的”这一事实。



 
对 (\ref{elx}) 的一个解释如下。按定义，$\J_I[\vec z_i] $ 是 $\U$ 中可由如下流
\begin{equation}
\frac{\d \bar w_{\bar i}}{\d t} = -e^{i \alpha} g^{\bar i j}  \frac{\partial \W(w,z_i)}{\partial w_j}
\end{equation}
到达的点的集合，这些流在过去渐近于 $I$。这里 $z_i$ 被看作常数。现在考虑半无穷区间 $(-\infty,t_0]$ 上受迫梯度流的方程 (\ref{flow})。对 $t_0\leq -T$，该方程的解在 $t_0$ 处的取值参数化了 $\J_I[\vec z_i]$；但对 $t_0>-T$，它们参数化了这个空间的一个依赖于 $t_0$ 的连续形变，我们将它记作 $\tilde \J_I[\vec z_i;t_0]$。

我们在第~\ref{dual}~节中看到，对 $H_m(\U ,\U_<)_{\vec z_f}$ 中任一给定的链 $\Gamma$，展开式
\begin{equation}
\Gamma =\sum_J\frak{c}^J \J_J[\vec z_f]
\end{equation}
中的系数 $\frak{c}^J $ 按指标的意义计数了如下流的数目
\begin{equation}
\frac{d \bar w_{\bar i}}{dt} = -e^{i \alpha} g^{\bar i j}  \frac{\partial \W(w,z_i)}{\partial w_j}
\end{equation}
这些流从 $\Gamma$ 出发，在未来渐近于临界点 $J$。

因此，在编织发生之后的某个给定时刻 $t_0$，从 $\tilde \J_I[\vec z_i;t_0]$ 流向未来临界点 $J$ 的四维解的数目，就是展开式
\begin{equation} \label{ntop}
\tilde \J_I[\vec z_i;t_0] =\sum_J\tilde {\frak N}_I^J\J_J[z_f].
\end{equation}
中的系数 $\tilde{\frak N}^J_I$。但由于 $\tilde\J_I[\vec z_i;t_0]$ 参数化了区间 $(-\infty,t_0]$ 上始于 $I$ 的流，从 $\tilde\J_I[\vec z_i;t_0]$ 到 $J$ 在区间 $[t_0,\infty)$ 上的流，等价于定义在整条实直线上的从 $I$ 到 $J$ 的流。所以 $\tilde{\frak N}^J_I$ 与所要的不变量 $\frak N^J_I$ 相同：
\begin{equation} \label{nzop}
\tilde \J_I[\vec z_i;t_0] =\sum_J\frak N^J_I  \J_J[z_f]
\end{equation}
相对同调定义在 $\Z$ 上，而像 $\tilde\J_I$ 这样的整相对同调类没有连续形变。因此很清楚：只要由流方程 (\ref{flow}) 诱导的、从 $\J_I[\vec z_i]$ 到 $\tilde \J_I[\vec z;t_0]$ 的连续形变在任一给定时刻都落在 $H_m(\U ,\U_<)_{\vec z(t)}$ 中，它就与沿 $\P$ 由 Gauss-Manin 联络给出的自然搬运一致。在这种情形下，(\ref{nzop}) 等价于所要的结果 $\frak{N}^J_I = \frak{n}^J_I$。

为证明对任意 $t_0$ 都有 $\tilde\J_I[\vec z;t_0]$ 落在 $H_m(\U,\U_<)_{\vec z(t_0)}$ 中，我们应当证明 $\mathrm{Re}\,\W(w,\vec z(t_0))$ 沿 $\tilde \J_I[t_0]$ 到无穷远趋于 $-\infty$。确实，若半无穷区间 $(-\infty,t_0]$ 上的一列受迫梯度流趋于无穷远，它一定是通过 $t\to t_0$ 时发散而趋于无穷远的；在这种情形下，$\mathrm{Re}(\W(w,z(t_0))$（其梯度驱动 $t\to t_0$ 时的流）必须趋于 $-\infty$。

对 (\ref{elx}) 的另一种论证如下。假定 $\vec z_f=\vec z_i$，且 $\P$ 是它们之间的平凡路径。那么 $\frak N^J_I=\frak n^J_I=\delta^J_I$。当我们变动 $\vec z_f$ 时，$\frak N^J_I$ 与 $\frak n^J_I$ 都只可能在穿过 Stokes 墙时跳跃；它们的跳跃方式完全相同，因此始终保持相等。我们已经描述了 $\frak n^J_I$ 的跳跃。$\frak N^J_I$ 的跳跃之所以发生，是因为在穿过 Stokes 墙时，一个依赖时间的解可能消失到无穷远，具体如下。假定对某个 Stokes 墙上的 $\vec z_f$，$\frak n^J_I$ 发生跳跃，其来源是从 $\W(w,\vec z_f)$ 的某个临界点 $J'$ 到 $J$ 的一条普通梯度流。这样的流产生跳跃
\be\label{jumpo} \frak n^J_I\to \frak n^J_I\pm \frak n^{J'}_I,\ee
其中符号取决于穿过 Stokes 墙的方向。为看到 $\frak N^J_I$ 中相应的跳跃，人们寻找从 $I$ 到 $J$ 的受迫梯度轨迹，它由一条从 $I$ 到 $J'$ 的受迫梯度轨迹、随后在某个非常遥远的未来时刻接上同一条引起 $\frak n^J_I$ 跳跃的从 $J'$ 到 $J$ 的普通梯度轨迹所组成。若 $\vec z_f$ 靠近 Stokes 墙且位于其正确的一侧，这种两步受迫轨迹就存在；当 $\vec z_f$ 穿过 Stokes 墙时，流发生第二步的时刻发散。这导致两步解消失，并使 $\frak N^J_I$ 产生一个恰好与 $\frak n^J_I$ 的跳跃相配的跳跃。




\section{从针的编织到自由场积分}\label{fthimconformal}
\begin{keytakeaways}
\begin{itemize}\setlength{\itemsep}{2pt}
  \item \textbf{物理图像}：把相对同调的链看作积分链，BPS 态随参数的演化就等价于这些积分随参数的演化；针上的积分在 $\varepsilon\to0$ 的半经典极限下由临界点主导，故可用最速下降/Picard-Lefschetz 方法来计算。
  \item \textbf{数学工具}：相对同调上的 Gauss-Manin 联络与 Stokes 理论、针上的渐近展开、复结构形变（Beltrami 微分）、Chern-Simons 路径积分与共形块。
  \item \textbf{主要结论}：$\Mthree=C\times\R_+$ 上针的积分等同于带 Nahm 极点边界条件的 $\N=4$ 理论在 $M_3\times\tilde\R_+$ 上的配分函数；两个 Chern-Simons 描述由 $S$ 对偶相联系，其能级满足 $\k+2=1/(\k^\vee+2)$。
  \item \textbf{共形块}：oper 边界条件（$\A_z$ 上三角且右上元为 1、$\A_{\bar z}$ 下三角）诱导出 Virasoro 共形块，且 Jones 参数与 $b$ 满足 $q=\exp(-2\pi i/b^2)$。
\end{itemize}
\end{keytakeaways}

根据第~\ref{fourd}~节的论证，为理解与 Jones 多项式相联系的辫群表示，我们应当计算三流形 $\Mthree=C\times \R_+$ 上联络的无穷维空间 $\U$ 的中间维相对同调上的一个自然单值性作用。这听起来可能抽象得令人绝望。现在我们将说明，它可以转化为某种具体且可计算的东西。


\subsection{从针到积分}\label{fromth}



描述 $\H$ 中态的演化（包括 $q$-分次）的一种方便方式，是把同调链看作积分链。不考察同调链的演化，等价地可以考察积分的演化。当然，我们需要一个可以在针 $\J_I$ 上积分的被积函数，而针不是紧的。满足这一要求的函数是 $e^{\W/\varepsilon}$，其中 $\varepsilon$ 选取在适当的半平面内。由于
\begin{equation}
h = \mathrm{Re} \left( e^{i \alpha} \W \right)
\end{equation}
沿针趋于负无穷，我们所要的条件是 $\mathrm{Re} \left( e^{i \alpha} \varepsilon\right)>0$。这将保证积分
\begin{equation}\label{thelko}
\I_\Gamma = \oint_\Gamma e^{\W(w)/\varepsilon} \d \Omega
\end{equation}
收敛，其中 $\d\Omega$ 是 $\U$ 上的全纯体积形式（下文将保持固定），而积分链 $\Gamma$ 是一个针，或更一般地是相对同调 $H(\U,\U_<)$ 中的任一链。

针是特别好的积分链，因为针上积分的 $\varepsilon\to 0$ 极限非常简单。在由从临界点 $I$ 出发的梯度流定义的针 $\J_I$ 上，函数 $h$ 有唯一的最大值，即该临界点。因此当 $\varepsilon\to 0$ 时，针上的积分为
\begin{equation}\label{thelk}
\I_I:= \oint_{\J_I} e^{ \W(w)/\varepsilon}\, \d \Omega\sim \exp(\W_I/\varepsilon)\left(\varepsilon^{-\mathrm{dim}\,\J_I/2}c_0+\dots\right),
\end{equation}
其中 $\W_I$ 是 $\W$ 在临界点 $I$ 处的值。
% [Tutor Note: 验证计算] (\ref{thelk}) 的幂次 $\varepsilon^{-\mathrm{dim}\,\J_I/2}$ 来自针上的高斯型积分：$\J_I$ 的实维数等于 $\U$ 的复维数（也就是这里所说的“中间维”），沿其做 $\mathrm{dim}\,\J_I$ 个实方向的鞍点积分，每个实方向贡献 $\varepsilon^{-1/2}$，合起来便是 $\varepsilon^{-\mathrm{dim}\,\J_I/2}$。注意不要把 $\mathrm{dim}\,\J_I$ 与某个复维数混用，否则幂次会差因子 $2$。

这个公式在整个半平面 $\mathrm{Re} \left( e^{i \alpha} \varepsilon\right)>0$ 内成立；但实际上只要 $\alpha$ 不在 Stokes 墙上，这个渐近式在复平面内稍大一些的扇区内也成立。\footnote{\label{breakdown}为使所述的渐近式失效，第一步是穿过一面 Stokes 墙，于是我们出发时的针演化为至少含两项的针的线性组合。起初，渐近式 (\ref{thelk}) 仍然成立，因为 Stokes 墙上新出现的那些额外的针起初只贡献指数小的量。若把 $\alpha$ 继续变动，其中一个额外针最终可能占主导。整个过程总是涉及把 $\alpha$ 从其初值变动一个严格大于 $\pi/2$ 的角度。为说明这一点，只需比较 $\W$ 在两个临界点处的值；这些值在 Stokes 墙上虚部相等，而在两个临界点交换主导地位时实部相等。}这一性质唯一地刻画了所有可能的 $I_\Gamma$ 中积分基 $\I_I$：若取几个 $\I_I$ 的线性组合，则渐近式会在扩展半平面内的某条射线上失效，在彼处两个临界点交换主导地位。这一刻画在 Stokes 理论中为人熟知，也正是“Stokes 墙”这一术语的由来。

当 $\W$ 的参数变动时，同调 $H(\U,\U_<)$ 连续地变动。若把积分链 $\Gamma$ 连续地变动，积分将全纯地依赖于 $\W$ 的参数。链 $\Gamma$ 的单值性与积分 $\I_\Gamma$ 的单值性相同。

\subsection{两个 Chern-Simons 理论}
在我们当前的语境中，针的积分 (\ref{thelko}) 是三流形 $\Mthree = C \times \R_+$ 上的 Chern-Simons 路径积分，只不过是在一个不寻常的积分链上。这样的积分可以用 $\Mthree \times \tilde \R_+$ 上的 $\N=4$ 规范场论具体地表达出来。（仅在本节中，我们把 $x^1$ 方向写作 $\tilde \R$ 或 $\tilde\R_+$，以区别于 $y$ 方向 $\R_+$。）这一陈述是文献~\cite{Analytic, NewLook} 的主要结论之一。在我们的情形中，由于 $\Mthree = C \times \R_+$，四维流形是 $M_4 = C \times \R_+ \times \tilde \R_+$。

这是本文中出现的第二个 Chern-Simons 理论，至少是隐含地出现。我们的全部分析都关乎 Jones 多项式的计算，所用的规范场论设定 $S$ 对偶于一个通过在 $W = C \times \tilde \R$ 上的 Chern-Simons 理论计算 Jones 多项式的设定。在那个“原始的” Chern-Simons 理论中，纽结是一条 Wilson 圈，规范群是 $G^\vee$，耦合参数是 $\k^\vee$。Jones 多项式是如下量的 Laurent 多项式：
\begin{equation}\label{torg} q=\exp(2\pi i/(\k^\vee+2)),\end{equation}
其中 $2$ 是 $SU(2)$ 的对偶 Coxeter 数。

正如文献~\cite{fiveknots} 所解释的，若把 $\k^\vee$ 与 $\N=4$ 理论的扭曲参数 $\Psi^\vee$ 通过
\begin{equation} \label{psikvee} \Psi^\vee=\pm (\k^\vee+2). \end{equation}
相联系，则三流形 $W$ 上的 Chern-Simons 理论可以通过 $W\times \R_+$ 上的拓扑扭曲 $\N=4$ 超 Yang-Mills 理论来计算。在这一描述中，人们在 $\R_+$ 的原点处使用 D3-NS5 边界条件。符号 $\pm$ 取决于 $M_4$ 与 $W$ 之间定向的相对选取。

人们还可以应用 $S$ 对偶，把规范群从 $G^\vee=SU(2)$ 变为 $G=SO(3)$，并把 D3-NS5 边界条件变成 D3-D5 边界条件；如文献~\cite{gw} 所解释的，这个边界条件涉及 Nahm 极点。在这个新的描述中（它正是本文的出发点），扭曲参数为 $\Psi=-1/\Psi^\vee$。由于这个对偶描述很困难，我们尝试按第~\ref{methods}~节首先说明的办法，通过沿一个方向“拉伸” $W$ 来简化它。于是我们用 $\tilde \R\times C$ 来近似 $W$，并寻找 $W\times \R_+$ 上从 $M_3=C\times \R_+$ “拉回”的 BPS 方程的解。若把这个对偶描述表述为以 $\R$ 为时间方向、以 $M_3=C\times \tilde \R_+$ 上的规范场为场变量的超对称量子力学，则超势是 $\Psi$（而非 $\Psi^\vee$）乘以 Chern-Simons 泛函。

然后我们把 $M_3$ 上解的编织与 $M_3$ 上 Chern-Simons 理论的针的编织联系起来，后者我们用 $M_3\times \tilde \R_+$ 上的 $\N=4$ 理论来表达。在这一描述中，我们在 $\tilde \R_+$ 的原点处施加 D3-NS5 边界条件（并在 $\R_+$ 的原点处施加 D3-D5 边界条件）。这一 Chern-Simons 描述的能级 $\k$ 与 $\N=4$ 理论的扭曲参数通过
\begin{equation} \label{psik} \Psi=\mp (\k+2). \end{equation}
相联系。这两个 Chern-Simons 描述由 S 对偶以及 $\R_+$ 与 $\tilde \R_+$ 的角色互换相联系。(\ref{psikvee}) 与 (\ref{psik}) 中符号相反，是因为交换 $\R_+$ 与 $\tilde \R_+$ 会反转 $M_4$ 的定向。
% [Tutor Note: 易错点] 两个符号 $\pm$ 与 $\mp$ 必须配套使用：$S$ 对偶同时把 $\Psi\to-1/\Psi^\vee$ 并交换 $\R_+\leftrightarrow\tilde\R_+$，后者反转 $M_4$ 定向，才导致 (\ref{psik}) 相对 (\ref{psikvee}) 取相反的符号。若忘记定向反转，会得到 $\Psi=+\ (\k+2)$，进而把 (\ref{sdualk}) 的倒数关系写成非倒数关系。

结合 $\Psi=-1/\Psi^\vee$ 与 $\Psi^\vee
=\k^\vee+2$，我们发现：与传统方式下联系 Jones 多项式的 Chern-Simons 描述中的能级 $\k^\vee$，和把 Jones 多项式与 Nahm 极点及 oper 相联系的 Chern-Simons 描述中的能级 $\k$，二者之间的关系是\begin{equation}\label{sdualk} \k+2=1/(\k^\vee+2).\end{equation}
% [Tutor Note: 物理图像] (\ref{sdualk}) 并不是简单的 $\k\leftrightarrow\k^\vee$：注意移位后的能级 $\k+2$ 才参与倒数关系，且两个理论互为强弱耦合（$\k^\vee\to\infty$ 时 $\k\to-2$）。$SU(2)$ 的对偶 Coxeter 数为 $2$，正是 (\ref{torg}) 中出现 $+2$ 的来源；这个 $2$ 在一般群 $G$ 时应换成对偶 Coxeter 数。

这个四维设定——$\N=4$ 理论处在带一个“角”的流形 $M_4 = C \times \R_+ \times \tilde \R_+$ 上——相当有趣，我们相信它值得进一步探讨。本文不再深入。

\subsection{从 Chern-Simons 到共形块}

对本文而言，比 $\Mthree=C \times \R_+$ 上 Chern-Simons 理论与四维的联系更有用的，是它与 $C$ 上共形块的联系。这一陈述最熟悉的版本（文献~\cite{WittenJones}）是：固定 $\A_{\bar z}$ 而让 $\A_z$ 变化的路径积分给出一个 WZW 共形块（即二维流代数中的共形块，其对称群 $G$ 与 Chern-Simons 理论的规范群相同）。对固定的 $\A_{\bar z}$ 作局部变动，就在共形块中插入一个全纯流 $J(z)$：
\begin{equation}
\delta \log \cZ = \int \mathrm{Tr} \langle J_z \rangle \delta A_{\bar z}
\end{equation}

这一陈述有一个类比（文献~\cite{HVerlinde}），在 $SU(2)$ 或 $SO(3)$ 规范场论的情形给出 Virasoro 共形块，而在更高秩规范群的情形给出更一般的 $W$-代数共形块（文献~\cite{Bilal,DeBoer}）。这个类比涉及一种不同的边界条件：以在边界上破坏一部分规范对称为代价，固定 $\A_z$ 的一部分和 $A_{\bar z}$ 的一部分。对规范群 $SU(2)$，导致 Virasoro 共形块的边界条件可简单地表述为：
\begin{align}\label{bcon}
\A_z = \begin{pmatrix} * & 1 \cr * & *\end{pmatrix} \qquad \A_{\bar z} = \begin{pmatrix} \times & 0 \cr * & \times\end{pmatrix}
\end{align}
这里我们把可以自由变动的元素记作 $*$，把固定的元素记作 $\times$。

刚才描述的联络就是一个 oper！（由于 $\A_{\bar z}$ 是下三角形的，带有这个联络的丛有一个全纯子丛 $L$，其截面形如 $\begin{pmatrix}0\cr *\end{pmatrix}$；又因为 $\A_z$ 的右上矩阵元处处不为零，这个子丛处处不被 $\D_z=\partial_z+[\A_z,\,\cdot\,]$ 保持。）所以复边界条件 (\ref{bcon}) 正是第~\ref{analog}~节所研究的 Nahm 极点边界条件所诱导的那个。为说明与 Virasoro 共形块的关系，我们指出：Nahm 极点边界条件依赖于 $C$ 上复结构的一个选取。一旦选定复结构并取局部复坐标 $z$，附近的复结构就可以用 Beltrami 微分 $\mu_{\bar z}^z$ 描述。其关系是：在新的复结构下，复联络空间上的全纯场不是 $\A_{\bar z}$，而是 $\A_{\bar z}-\mu_{\bar z}^z\,\A_z$。作这样的形变，等价于把 (\ref{bcon}) 中 $\A_{\bar z}$ 边界条件里的“$0$”换成 $\mu_{\bar z}^z$（于是现在 $\A_{\bar z}-\mu_{\bar z}^z\A_z$ 是下三角形的）。因此，与复结构改变相联系的边界条件的局部变动会插入一个全纯应力张量
\begin{equation}
\delta \log \cZ = \int \langle T_{zz} \rangle \delta \mu_{\bar z}^z,
\end{equation}
而这导致带有 oper 边界条件的 Chern-Simons 理论与 Virasoro 共形块之间的关系。
% [Tutor Note: 物理图像] 这段把“oper 边界条件”与“复结构形变”联系起来：$\A_z$ 的右上元固定为 1 使得子丛 $L$（截面 $(0,*)$）不被 $\D_z$ 保持（oper 的“不可积”条件），而 $\A_{\bar z}$ 的下三角形式让 $L$ 全纯；Beltrami 微分 $\mu_{\bar z}^z$ 描述复结构的形变，其系数正是应力张量 $T_{zz}$ 的源，于是对边界条件求变分就得到 Virasoro Ward 恒等式的插入。

在半经典极限下，算符 $T_{zz}$ 化为 oper 的经典应力张量 $\bt(z)[\A]$。更准确地说，文献~\cite{HVerlinde} 的参数对应关系是：若定义
\begin{equation}\label{bpar}-b^{-2} =\k+2,\end{equation}
其中 $\k$ 是 Chern-Simons 能级，则应力张量的中心荷为 $c=1+6(b+b^{-1})^2$，且在半经典极限下 $b^2 T(z) \to \bt(z)[\A]$。

若把 (\ref{bpar}) 与 (\ref{sdualk})、(\ref{torg}) 结合起来，我们就得到通常用于描述 Jones 多项式的变量 $q$ 与用于描述 Virasoro 共形块的参数 $b$ 之间的关系：
\begin{equation}\label{zorg}q=\exp(-2\pi i/b^2).\end{equation}

这里我们假定了 $G$ 是 $SU(2)$ 或 $SO(3)$。对一般的规范群 $G$，Nahm 极点与 $W$-代数都由 $G$ 的李代数 $\frak{g}$ 中的一个 $\frak{su}(2)$ 嵌入来标记。检查可以确认，用于定义一般 $W$-代数共形块的边界条件正是由相应的 Nahm 极点所诱导的。
